\section{The cost of partial transposition}\label{sec:partial}

Transposing a binary matrix makes each input bit available at a new
place. Bit $(i,k)$ enters with row $i$ and leaves with column $k$.
On tapes the heads can decide where to move from the data they have
already seen. Fixing their movements at a coarse resolution gives a
small network, but only for the inputs that follow those movements.
We first bound the size of a network that transposes correctly on
even an exponentially small set of inputs.

\begin{theorem}[Partial transposition]\label{thm:table}
Assume Weak ETH and fix $0\le\rho<1-1/\sqrt2$. Let a rational tensor
network have maximum degree $D$ in edge bundles, each of width at most
$DN$ bits. Give its $N$ input row strings and $N$ output column strings
of length $N$ distinct terminals. Suppose that on a fraction at least
$2^{-\rho N^2}$ of all binary $N\times N$ inputs, fixing the row strings
gives a nonzero scalar multiple of the basis vector for the transposed
output. Then the network has $\Omega_{D,\rho}(N\sqrt{\log N})$ vertices
for every sufficiently large $N$.
\end{theorem}

Two estimates give the bound. The first places the network so that the
cost of its edges falls by a factor of order $\sqrt{\log N}$ compared
with that of typical row-to-column wires. The second turns the entropy
of the successful inputs into a lower bound on the cost of those wires.

\subsection{Charging the wires of a small network}

Represent a shared string of $q$ bits by a wire of width $q$. Place the
vertices of the network in a spectral space from Section~\ref{sec:graphs}.
The price of each wire is its width times the distance between its ends.
If $q_e$ is the width of edge $e$, the total supply price is
\begin{equation}\label{eq:supply-cost}
 S=\sum_{e=uv}q_e\,d(x_u,x_v).
\end{equation}

We need a placement with three properties. The network edges must be
short, typical row and column terminals must be far apart, and there
must be a central point from which we can identify the chosen input
rows at moderate cost. We normalize the typical terminal distance to one.

\begin{lemma}[Short edges and distant terminals]\label{lem:geometry}
Assume Weak ETH. Let a multigraph have $P_0$ vertices, maximum
degree $D$, and edge widths at most $Db$. For all sufficiently large $P_0$
there is a placement in some $X\in\mathcal X$ and a point $o\in X$ with
the following properties. Write $P$ for the least power of two at least
$P_0$. All constants below depend only on $D$.
\begin{align}
 S &\le C_D bP/\sqrt{\log P},                                           \label{eq:short-supply}\\
 \E_{a\in A,\,t\in B}\bigl(d(x_a,x_t)-1\bigr)^2
   &\le C_D P/(N\log P),                                                \label{eq:terminal-distance}\\
 \sum_v h_v d(x_v,o)
   &\le \frac1{\sqrt2}\sum_v h_v+
       C_D\left(\frac{P}{\log P}\sum_v h_v^2\right)^{1/2}.             \label{eq:central-cost}
\end{align}
Here $A,B$ are any two lists of $N$ vertices distinct within each list,
the draws from the two lists are independent and uniform, and $h_v\ge0$
are arbitrary. In particular, when $P=O(N\sqrt{\log N})$,
the mean absolute deviation in \eqref{eq:terminal-distance} is $o(1)$.
If also $\max h_v=O(b)$ and $\sum h_v=O(Nb)$, the error in
\eqref{eq:central-cost} is $o(Nb)$. These bounds are uniform over the lists, and the weights with the indicated bounds.
\end{lemma}

\begin{proof}
Add vertices to reach $P$ and superimpose a scalar expander on all $P$
vertices, of fixed degree and with a fixed spectral gap \cite{HLW,Friedman}.
Add diagonal arcs to make every degree equal to a fixed power of two
$d_0\ge2$ depending only on $D$. Each nonloop edge gives two reverse arcs, and each diagonal arc
counts once.
Under Weak ETH (Theorem~\ref{thm:noports}), the logarithmic obstruction
\cite[Theorem~\ext{P}{thm:loggraph}]{oaP}, based on
nonlinear spectral calculus \cite{MN}, applies to $\mathcal X$
(a class of Hadamard spaces without ports).
It says that the best uniform Poincar\'e constant of
this graph over $\mathcal X$ (in the squared form)
is at least
$c_{d_0}\log P$. Thus for $P$ large some placement in a rescaled space has
\[
 \E d(x_U,x_V)^2=1,\qquad
 \E d(x_I,x_J)^2=O_D(1/\log P).
\]
Here $U,V$ are independent uniform vertices and $(I,J)$ a uniform
directed arc of the augmented graph. Its $d_0P$ arcs include the
original edges. Multiplying their mean length by $d_0P$ and by the
maximum width $Db$ gives \eqref{eq:short-supply} by Cauchy--Schwarz.
The scalar Poincar\'e
inequality on the expander subgraph now gives, for every real
$1$-Lipschitz $f$ on $X$,
\begin{equation}\label{eq:scalar-variance}
 \operatorname{Var} f(x_U)\le C_D/\log P.
\end{equation}
In particular, the bound holds for distance to any fixed point and for
$g(x)=\E d(x,x_U)$. The conditional variance formula, applied first with
$V$ fixed and then to $g(x_V)$, gives
$\operatorname{Var}d(x_U,x_V)=O_D(1/\log P)$. Since the second moment
is one, the mean $\mu$ is $1+O_D(1/\log P)$.

A list of $N$ distinct vertices has uniform density $N/P$.
Consequently its average of a nonnegative function is at most $P/N$
times the average over all vertices. For every $t$ we therefore have
$\E_{a\in A}(d(x_a,x_t)-g(x_t))^2\le C_D P/(N\log P)$.
The same bound holds for $\E_{t\in B}(g(x_t)-\mu)^2$.
Adding these deviations in $L^2$ proves \eqref{eq:terminal-distance}.

Take $o$ to be the barycenter of the uniform distribution on the $P$
points. In a complete $\operatorname{CAT}(0)$ space the mean squared
distance to a finite distribution has a unique minimizer $o$, and its
value at $y$ exceeds its minimum by at least $d(o,y)^2$.
Both facts follow from strong convexity along geodesics, first on
midpoints of a minimizing sequence and then along the segment from
$o$ to $y$ \cite{BH,Sturm}. Averaging the inequality at $y=x_V$ gives
$2\E d(x_U,o)^2\le1$. Apply \eqref{eq:scalar-variance} to $d(\,\cdot\,,o)$
and extend the weights by zero to the added vertices. Cauchy--Schwarz
bounds the weighted deviations from the mean by the last term of
\eqref{eq:central-cost}.
\end{proof}

\subsection{Turning successful inputs into independent wires}

On a good input each bit of an array appears at two prescribed
terminals. The bits need not be independent, and the network may
produce an incorrect output on other inputs. To charge a price for
the bits that do appear, allow a central source at $o$ to select the
local tables for a good input. A \emph{pin} is the index of one such
choice at a vertex, incorporating all choices at that vertex. 
Lemma~\ref{lem:entropy} charges for transmitting the pins and credits
each terminal pair with the new information in its bit. Write $H$ for
Shannon entropy in bits.

\begin{lemma}[Entropy at paired terminals]\label{lem:entropy}
Let a network have pairs of indicated vertices (terminals)~$(a_j,t_j)$, $1\le j\le L$, and let
$W=(W_1,\ldots,W_L)$ be a random binary vector. For each $w$ in its support,
suppose pins $C_v(w)$ at the vertices make the contracted tensor
over the rationals equal to a nonzero scalar multiple of
\[
 \bigotimes_{j=1}^L\bigl((e_{w_j})_{a_j}\otimes(e_{w_j})_{t_j}\bigr).
\]
For every placement in $X\in\mathcal X$, every $o\in X$, and every
ordering of the coordinates, numbered in that order, we have
\begin{equation}\label{eq:entropy-price}
 \sum_{j=1}^{L} H(W_j\mid W_1,\ldots,W_{j-1})\,d(x_{a_j},x_{t_j})
 \ \le\ S+\sum_v H(C_v(W))\,d(o,x_v).
\end{equation}
\end{lemma}

\begin{proof}
To apply the price comparison we must extract a product of independent
equality wires. Over many uses the two endpoints of pair $j$ see the
same binary column. We assign a label to that column so that the labels
can be chosen independently, at a rate approaching the conditional
entropy of coordinate $j$.

It suffices by continuity to treat a rational probability law $p$ on
the given support. Take the tensor product of $r$ copies of the network, allowing joint
local maps at each vertex, and take $r$ through multiples of a common
denominator of $p$. We use arrays of $r$ payloads having joint type $p$,
meaning that each $w$ occurs exactly $rp(w)$ times. The number of
possible $j$th columns compatible with a fixed prefix of the prescribed
type is
\[
 M_j=\prod_{u\in\{0,1\}^{j-1}}
       \binom{r\,p(W_{<j}=u)}{r\,p(W_{<j}=u, W_j=1)}.
\]
Each such extension still has the prescribed prefix type and can be
continued to a full array. The type-counting estimate gives
$r^{-1}\log M_j\to H(W_j\mid W_{<j})$ \cite[Theorem~11.1.3]{CT}.

We need a map of the $j$th column alone, since its endpoints see only
that column. Put $\ell_j=\max(0,\lfloor\log(M_j/r^2)\rfloor)$.
There is a coloring of the binary $r$-columns by $2^{\ell_j}$ colors
such that every compatible prefix admits every color as an extension.
Indeed, if $\ell_j>0$, independent uniform colors miss a given color
among the $M_j$ extensions with probability at most $\exp(-r^2)$.
There are at most $2^{r(j-1)}$ prefixes and $2^r$ colors, so the union
bound applies for large $r$ with $L$ fixed. For $\ell_j=0$ the claim
is immediate. Choosing columns successively, every tuple of colors
has a preimage array of type $p$. Fix one such array for each tuple.

We now supply exactly the arrays selected by these labels. At $v$
each array determines a stream of $r$ pins whose type is the law
of $C_v(W)$. There are at most $2^{rH(C_v(W))}$ such streams, so a wire
of $\lceil rH(C_v(W))\rceil$ bits from a new vertex at $o$ can supply
its index. A basis index for a stream $(c_1,\ldots,c_r)$ selects at $v$
the tensor product of its $r$ prescribed local tables. This defines a
linear local map by extension from the basis indices. At $o$ form the sum, over the color tuples, of the basis vectors for
their joint pin indices. Distinct selected arrays give distinct pin
streams somewhere in the network, because the pins determine every
displayed bit. Each selected array has the same output scalar
$c_*=\prod_w c(w)^{rp(w)}$, where $c(w)\ne0$ is the one-copy scalar.
Divide the tensor at $o$ by $c_*$. Applying the coloring map at both
ends of every column now gives each possible tuple of labels exactly
once, with coefficient one. The output is therefore the product of
$\ell_j$-bit equality wires between $a_j$ and $t_j$.

The original edges have total cost $rS$. Including the wires from $o$,
the cost comparison using \eqref{eq:supply-cost} reads
\[
  \sum_j \ell_j d(x_{a_j},x_{t_j})
     \le rS+\sum_v \lceil rH(C_v(W))\rceil\,d(o,x_v).
\]
Divide by $r$ and let $r$ tend to infinity.
\end{proof}



\subsection{The price gap for partial transposition}

\begin{proof}[Proof of Theorem~\ref{thm:table}]
Write $n=N^2$ and $A_i,B_k$ for row and column terminals.
Take $W$ uniform on the good arrays. At terminal $A_i$, pin
its row and copy it to both the input and fresh external indices
there. The resulting tensor has a pair of external copies of every
bit of $W$, at its row and column (Figure~\ref{fig:advice}).
It suffices to consider networks with $P_0\le N\sqrt{\log N}$ vertices.

In the placement of Lemma~\ref{lem:geometry} with $b=N$ the error in pricing all
$N^2$ terminal-pair distances at one unit, in any sum with weights in $[0,1]$,
is $o(n)$. Thus the left side of \eqref{eq:entropy-price} is
$H(W)$ up to $o(n)$ by the chain rule. Since row entropies are at most
$N$, we can use \eqref{eq:central-cost} on the right side at a cost of
$\le n/\sqrt2+o(n)$ for delivering pins (recall all the relevant weights
are nonnegative and vanish off row terminals). We obtain
\[
 (1-\rho)n \ \le\
 O_D\!\left(\frac{N P_0}{\sqrt{\log P_0}}\right)+\frac n{\sqrt2}+o(n).
\]
The gap $1-\rho-1/\sqrt2$ is strictly positive, while $P_0\ge 2N$
by the distinct terminal assumption. The network must therefore have
$P_0=\Omega_{D,\rho}(N\sqrt{\log N})$ vertices.
\end{proof}

\begin{figure}[!htb]
\centering
\begin{tikzpicture}[font=\small,x=1cm,y=1cm,line cap=round,
 term/.style={minimum height=.43cm,minimum width=.52cm,inner sep=2pt,
              draw=ink,fill=white,rounded corners=.8pt},
 source/.style={draw=sea,fill=sea!5,circle,inner sep=3pt}]
\begin{scope}[xshift=-3.3cm]
  \node[font=\small\bfseries] at (0,1.55){Supplying the rows};
  \filldraw[fill=ink!4,draw=ink!35] (-.38,-1.08) rectangle (2.23,1.08);
  \node[align=center,text=ink,font=\footnotesize] at (.92,0)
     {network\\of price $S$};
  \node[source] (root) at (-2.28,0) {$o$};
  \foreach \i/\y in {1/.78,2/0,3/-.78}{
    \node[term] (A\i) at (-.38,\y) {$A_\i$};
    \node[term] at (2.23,\y) {$B_\i$};
    \draw[sea,line width=.6pt] (root)--(A\i);
  }
  \node[font=\footnotesize,align=center] at (.03,-1.53)
    {Pin cost at most $n/\sqrt2+o(n)$};
\end{scope}
\draw[ink,-{Stealth[length=2mm]}] (-.56,0)--(.56,0);
\begin{scope}[xshift=3.3cm]
  \node[font=\small\bfseries] at (0,1.55){Extracting independent wires};
  \foreach \i/\y in {1/.78,2/0,3/-.78}{
     \node[term] (left\i) at (-1.58,\y) {$A_\i$};
     \node[term] (right\i) at (1.58,\y) {$B_\i$};
  }
  \foreach \i in {1,...,3}\foreach\j in {1,...,3}
    \draw[ink!33,thin] (left\i.east)--(right\j.west);
  \draw[rust,line width=.9pt] (left2.east)--(right1.west);
  \node[font=\footnotesize,align=center] at (0,-1.53)
    {Wire cost at least $(1-\rho)n-o(n)$};
\end{scope}
\end{tikzpicture}
\caption{An array bit in position $(2,1)$ appears at row terminal $A_2$
and column terminal $B_1$. Over many uses, Lemma~\ref{lem:entropy}
turns these paired copies into independent wires, with a rate for
each pair given by its conditional entropy. The prices shown are per
use, with $n=N^2$ as $N$ grows. Three rows and columns illustrate the
pairing.}
\label{fig:advice}
\end{figure}


\subsection{From adaptive heads to a network}

\begin{theorem}[Transposition on a subset of inputs]\label{thm:transpose}
Assume Weak ETH and fix $0\le\rho<1-1/\sqrt2$. On a fixed deterministic
multitape machine $\mathcal M$, suppose that by time
$T$ it has halted with the correct transpose on a fraction at least
$2^{-\rho N^2}$ of all binary $N\times N$ matrices. Inputs and outputs occupy prescribed contiguous
intervals in row-major order. The initial machine configuration may
include auxiliary data fixed independently of the matrix. Then, at
every sufficiently large side length $N$,
\[
                 T=\Omega_{\mathcal M,\rho}(N^2\sqrt{\log N}).
\]
The constants and threshold are uniform over the auxiliary data.
In particular the bound holds for the worst-case time of an exact
transposition machine.
\end{theorem}

\begin{proof}
Fix the auxiliary data and let $\mathcal G$ be the matrices correctly
transposed by time $T$. We may assume $T\le N^2\sqrt{\log N}$. Divide
each tape into blocks of $N$ consecutive cells and the first $T$ steps
into epochs of $N$ steps. Use at least one epoch, pad the final epoch,
and freeze each halted configuration while enforcing the deadline.

At each epoch boundary record just the block containing each head.
A head's recorded block number changes by at most one per epoch, so
there are at most $2^{O_{\mathcal M}(1+T/N)}$ records with the fixed
starting positions. Keep a record shared by the largest number of
inputs in $\mathcal G$. The fraction of \emph{all} matrices in this
set is at least
\[
             2^{-\rho N^2-O_{\mathcal M}(1+T/N)}.
\]
Choose any fixed $\rho'$ between $\rho$ and $1-1/\sqrt2$. Since
$T/N=O(N\sqrt{\log N})=o(N^2)$, the fraction is at least
$2^{-\rho'N^2}$ for all sufficiently large $N$. Individual head
offsets and scanned symbols can still differ on these inputs.

Build a network with one vertex per epoch, using the recorded block
numbers to specify the graph (compare the blocking method of
\cite{hpv}). During an epoch a head can meet only its starting block
and its two neighbors. For each of these blocks carry its contents
after the epoch to the next epoch in whose neighborhood it appears.
Pass the control state and each head's offset in its boundary block
to the next epoch. The contents of a block require $O_{\mathcal M}(N)$
bits and all offsets and control information require
$O_{\mathcal M}(1+\log N)$ bits.

Figure~\ref{fig:epoch-block} shows how an epoch can pass the contents of
a block directly over an intervening epoch that cannot reach it.

\begin{figure}[!htb]
\centering
\begin{tikzpicture}[font=\small,
 epoch/.style={draw=ink,rounded corners=1.5pt,fill=ink!4,align=center,
               inner xsep=5pt,inner ysep=4pt,minimum width=2.54cm},
 state/.style={draw=ink,-{Stealth[length=1.6mm]},line width=.5pt},
 memory/.style={draw=sea,-{Stealth[length=1.7mm]},line width=.8pt}]
\foreach \i/\x/\b/\neighbors in
  {1/0/4/{3,4,5},2/3.45/5/{4,5,6},3/6.90/6/{5,6,7},4/10.35/5/{4,5,6}}{
  \node[epoch] (E\i) at (\x,0)
    {\textbf{Epoch \i}\\ start in block $\b$\\ $\{\neighbors\}$};
}
\foreach \i/\j in {1/2,2/3,3/4}
   \draw[state] (E\i.east)--(E\j.west);
\draw[memory] ([xshift=3mm]E1.south) -- ++(0,-.56)
   -| ([xshift=-3mm]E2.south);
\draw[memory] ([xshift=3mm]E2.south) -- ++(0,-.56)
   -| ([xshift=-3mm]E4.south);
\node[font=\footnotesize,text=sea,fill=white,inner sep=1.5pt] at (1.725,-1.16)
      {block $4$};
\node[font=\footnotesize,text=sea,fill=white,inner sep=1.5pt] at (6.9,-1.16)
      {block $4$};
\end{tikzpicture}
\caption{Four epochs on one tape, each labeled by the blocks it may reach.
The horizontal edges communicate the
controller state and head offset at successive boundaries. The lower
edges carry the contents of block $4$ between possible uses. A head
starting in block $6$ during Epoch 3 cannot reach block $4$ in $N$ steps,
so its contents pass directly from Epoch 2 to Epoch 4.}
\label{fig:epoch-block}
\end{figure}

The local table for an epoch is the indicator of a legal internal
run with the specified incoming and outgoing block contents, states
and head positions. Feed each block its initial contents at first use
and read its final contents at last use. Any unused block whose
initial contents affect the output can pass them through directly.
The final state must be halted by the deadline. Thus local runs
with matching boundary data concatenate to an actual machine run,
and every run that halts by the deadline with this record satisfies the
local relations. By
determinism, on each selected input the network has support only at
the correct output, with positive coefficient.

The fixed tape count bounds the number of incident bundles per epoch.
Use a distinct terminal for each row and column. An external string
crossing a block boundary is joined from at most two consecutive
blocks, and a block meets at most two such strings at either interface.
These extra vertices preserve bounded degree. In total there are
$O_{\mathcal M}(N+1+T/N)$ vertices, with all widths
$O_{\mathcal M}(N)$. Applying Theorem~\ref{thm:table} at density
$2^{-\rho'N^2}$ gives
\[
             N+1+T/N=\Omega_{\mathcal M,\rho}(N\sqrt{\log N}),
\]
as required.
\end{proof}

\begin{proof}[Proof of the multiplication lower bound in Theorem~\ref{thm:lower}]
Let $M(m)$ be the worst-case time of a fixed exact multiplier on two
$m$-bit operands. Given $m$, the reduction of Harvey and van der Hoeven
\cite[full version, Corollaries~6.2 and~6.3]{hvhtranspose} transposes
a bit matrix of side $N=\lfloor\sqrt m\rfloor$ in
\[
                O\bigl(M(m)+m(1+\log\log(2m))\bigr)
\]
steps. In that reduction, padding to a larger dyadic length is followed
by splitting each padded multiplication into a bounded number of calls
at length $m$. Thus the bound uses the actual value $M(m)$ at the
specified length. Copying, padding and adapting the ordinary binary
input and output conventions here take $O(m)$ additional steps.

Use $m$ as auxiliary data in Theorem~\ref{thm:transpose}, with
$\rho=0$. The lower bound on the resulting transposer is uniform in
$m$. Since $N^2\sqrt{\log N}=\Theta(m\sqrt{\log m})$ and
$m(1+\log\log(2m))=o(m\sqrt{\log m})$, it follows that
$M(m)=\Omega(m\sqrt{\log m})$ at every sufficiently large $m$.
\end{proof}
